% Master blueprint content document covering Phase 1 (F00-F06) and Layers 2-4 Roadmap.

\chapter{Foundations and Bookkeeping Layer}

\section{Toolchain and Dependency Audit}

\begin{definition}[Library and Toolchain Audit]
\label{f00:library_audit}
\lean{GromovCompactness.Audit.LibraryAuditRecord}
\leanok
Audit record certifying availability and compatibility of core external dependencies from Lean 4 (v4.16.0) and Mathlib4,
including Hofer's lemma (\texttt{Mathlib.Analysis.Hofer}), smooth manifolds with corners,
Riemannian metric structures, Hausdorff distance, Arzelà-Ascoli compactness, and acyclic graphs.
\end{definition}

\section{Target Geometry and the Cauchy-Riemann Equation}

\begin{definition}[Closed Tame Target]
\label{f01:target_geometry}
\lean{GromovCompactness.ClosedTameTarget}
\leanok
\uses{f00:library_audit}
Let $M$ be a compact, boundaryless, finite-dimensional smooth manifold. A \emph{closed tame target} consists of a closed non-degenerate differential two-form $\omega \in \Omega^2(M)$ and a smooth almost complex structure $J: TM \to TM$ ($J^2 = -\operatorname{id}_{TM}$) satisfying the tameness condition:
\[
\forall x \in M, \; \forall v \in T_x M \setminus \{0\}, \quad \omega_x(v, J_x v) > 0.
\]
\end{definition}

\begin{definition}[Symmetrized Riemannian Metric Tensor]
\label{def:symmetrized_metric}
\uses{f01:target_geometry}
\lean{GromovCompactness.ClosedTameTarget.metric}
\leanok
For a closed tame target $(M, \omega, J)$, the symmetrized Riemannian metric $g$ is defined on tangent vectors $v, w \in T_x M$ by:
\[
g_x(v, w) := \frac{1}{2}\big(\omega_x(v, J_x w) + \omega_x(w, J_x v)\big).
\]
\end{definition}

\begin{theorem}[Metric Properties: Symmetry, Positive Definiteness, and $J$-Invariance]
\label{thm:metric_pos_def}
\lean{GromovCompactness.metric_pos_def}
\leanok
\uses{f01:target_geometry, def:symmetrized_metric}
The symmetrized tensor $g$ is a genuine Riemannian metric on $M$:
\begin{enumerate}
\item \textbf{Symmetry}: $g(v, w) = g(w, v)$ for all $v, w \in TM$.
\item \textbf{Positive Definiteness}: $v \neq 0 \implies g(v, v) = \omega(v, Jv) > 0$.
\item \textbf{$J$-Invariance}: $g(Jv, Jw) = g(v, w)$ for all $v, w \in TM$.
\end{enumerate}
\end{theorem}
\begin{proof}
\leanok
Follows by algebraic symmetry of the polarizing sum, tameness of $J$, and $J^2 = -I$.
\end{proof}

\begin{definition}[Riemann Domain and Cauchy-Riemann Operator]
\label{def:j_holomorphic}
\lean{GromovCompactness.JHolomorphicOn}
\leanok
\uses{f01:target_geometry}
Let $(\Sigma, j)$ be a smooth Riemann surface domain. A smooth map $u : \Sigma \to M$ is called \emph{$J$-holomorphic} on a subset $s \subseteq \Sigma$ if it satisfies the Cauchy-Riemann equation:
\[
\delbar_J u := \frac{1}{2}(du + J \circ du \circ j) = 0 \iff du \circ j = J \circ du.
\]
\end{definition}

\begin{theorem}[Coordinate Invariance of the Cauchy-Riemann Operator]
\label{thm:jhol_coordinate_invariance}
\lean{GromovCompactness.jholomorphic_coordinate_invariance}
\leanok
\uses{def:j_holomorphic}
If $u : \Sigma \to M$ is $J$-holomorphic on $s$ and $\psi : \Sigma \to \Sigma$ is a biholomorphic diffeomorphism ($\psi^* j = j$), then $u \circ \psi$ is $J$-holomorphic on $\psi^{-1}(s)$.
\end{theorem}
\begin{proof}
\leanok
Direct computation via the chain rule on tangent bundles: $d(u \circ \psi) \circ j = du \circ d\psi \circ j = du \circ j \circ d\psi = J \circ du \circ d\psi = J \circ d(u \circ \psi)$.
\end{proof}

\section{The Riemann Sphere and Möbius Frame Transformations}

\begin{definition}[Riemann Sphere and Möbius Frames]
\label{f02:sphere_geometry}
\lean{GromovCompactness.MoebiusFrame}
\leanok
\uses{f00:library_audit, f01:target_geometry}
The standard Riemann sphere $\CP \cong S^2 \subset \R^3$ equipped with its canonical complex structure $j_{\CP}$.
The transformation group $\PSL$ acts transitively and biholomorphically as smooth Möbius frames, generated by affine scalings $z \mapsto a + \lambda z$, rotations $z \mapsto e^{i\theta} z$, and conformal inversion $z \mapsto 1/z$.
\end{definition}

\begin{theorem}[Three-Point Transitivity and Uniqueness in $\PSL$]
\label{thm:three_point_transitivity}
\lean{GromovCompactness.exists_unique_moebius_three_points}
\leanok
\uses{f02:sphere_geometry}
For any two triples of pairwise distinct points $(z_1, z_2, z_3)$ and $(w_1, w_2, w_3)$ in $\CP$, there exists a unique Möbius frame $\phi \in \PSL$ such that $\phi(z_i) = w_i$ for $i \in \{1, 2, 3\}$.
\end{theorem}
\begin{proof}
\leanok
Constructed via projective cross-ratio matrix coordinates in $\GL / \C^\times$.
\end{proof}

\section{Symplectic Energy Measure Algebra}

\begin{definition}[Symplectic Energy Measure]
\label{f03:energy_measure}
\lean{GromovCompactness.energyMeasure}
\leanok
\uses{f01:target_geometry, f02:sphere_geometry}
For a smooth map $u : \CP \to M$, the energy measure on a Borel subset $A \subseteq \CP$ is defined by:
\[
E(u; A) := \frac{1}{2} \int_A |du|_g^2 \, \dvol_{\CP}.
\]
The total energy is $E(u) := E(u; \CP) \in [0, \infty)$.
\end{definition}

\begin{theorem}[Energy Density Equals Pullback of $\omega$]
\label{thm:energy_eq_pullback_omega}
\lean{GromovCompactness.energy_eq_pullback_omega}
\leanok
\uses{f03:energy_measure}
For any $J$-holomorphic curve $u : \CP \to M$, the Riemannian energy density coincides pointwise with the pullback symplectic form:
\[
\frac{1}{2} |du|_g^2 \, \dvol_{\CP} = u^* \omega.
\]
Consequently, for any Borel set $A \subseteq \CP$, $E(u; A) = \int_A u^* \omega$.
\end{theorem}
\begin{proof}
\leanok
Follows from $J$-invariance of the metric $g$, tameness, and $du \circ j = J \circ du$.
\end{proof}

\begin{theorem}[Conformal Invariance of Symplectic Energy]
\label{thm:energy_conformal_invariant}
\lean{GromovCompactness.energy_conformal_invariant}
\leanok
\uses{f03:energy_measure, f02:sphere_geometry}
For any biholomorphic Möbius transformation $\phi \in \PSL$ and any Borel set $A \subseteq \CP$:
\[
E(u \circ \phi; \phi^{-1}(A)) = E(u; A).
\]
\end{theorem}
\begin{proof}
\leanok
By the change of variables formula for pullback differential forms under orientation-preserving conformal diffeomorphisms.
\end{proof}

\begin{theorem}[Energy Measure Null Sets]
\label{thm:energy_null_sets}
\lean{GromovCompactness.energy_point_null}
\leanok
\uses{f03:energy_measure}
For any point $p \in \CP$ and concentric round circle $C \subset \CP$:
\[
E(u; \{p\}) = 0 \quad \text{and} \quad E(u; C) = 0.
\]
\end{theorem}
\begin{proof}
\leanok
Singletons and smooth 1-manifolds have 2-dimensional Hausdorff measure zero.
\end{proof}

\section{Smooth Local Convergence}

\begin{definition}[$C^\infty_{\mathrm{loc}}$ Convergence on Exhausting Domains]
\label{f04:smooth_local_convergence}
\lean{GromovCompactness.SmoothLocalConvergesOn}
\leanok
\uses{f01:target_geometry, f02:sphere_geometry}
A sequence of smooth maps $u_n : \CP \to M$ converges locally smoothly in $C^\infty_{\mathrm{loc}}(\Omega, M)$ to $u : \Omega \to M$ on an open set $\Omega \subseteq \CP$ if for every compact subset $K \subset \Omega$, $u_n|_K \to u|_K$ uniformly with all partial derivatives of all orders.
\end{definition}

\begin{theorem}[Cauchy-Riemann Closure Under Smooth Local Limits]
\label{thm:smooth_local_jhol_closure}
\lean{GromovCompactness.smooth_local_jhol_closure}
\leanok
\uses{f04:smooth_local_convergence}
If $u_n \to u$ in $C^\infty_{\mathrm{loc}}(\Omega, M)$ and each $u_n$ is $J$-holomorphic on $\Omega$, then the limit map $u$ is $J$-holomorphic on $\Omega$.
\end{theorem}
\begin{proof}
\leanok
The differential operator $\delbar_J$ is continuous with respect to the $C^1_{\mathrm{loc}}$ topology, hence $\delbar_J u = \lim_{n \to \infty} \delbar_J u_n = 0$.
\end{proof}

\section{Subsequence Extraction and Diagonalization}

\begin{theorem}[Countable Diagonal Subsequence Extraction]
\label{f05:subsequence_bookkeeping}
\lean{GromovCompactness.Bookkeeping.countable_diagonal_extraction}
\leanok
\uses{f00:library_audit}
Given a countable sequence of nested strictly increasing subsequences $\sigma_{k+1} \prec \sigma_k$ (each refining the previous), there exists a diagonal subsequence $\sigma_{\mathrm{diag}} : \N \to \N$ that is strictly monotone and eventually a subsequence of every level $\sigma_k$.
\end{theorem}
\begin{proof}
\leanok
Formalized via Cantor diagonal extraction $\sigma_{\mathrm{diag}}(n) = \sigma_n(n)$ with tail bounds.
\end{proof}

\section{Combinatorial Bubble Trees and Stability}

\begin{definition}[Stable Sphere Tree Limit]
\label{f06:sphere_tree}
\lean{GromovCompactness.StableSphereTreeLimit}
\leanok
\uses{f01:target_geometry, f02:sphere_geometry, f03:energy_measure, f04:smooth_local_convergence, f05:subsequence_bookkeeping}
A finite acyclic connected tree $T = (V, E_T)$ equipped with:
\begin{enumerate}
\item Component curves $u_v : \CP \to M$ for each vertex $v \in V$,
\item Distinct nodal attachment points $z_{vw} \in \CP$ for each directed edge $(v, w) \in E_T$,
\item Nodal matching $u_v(z_{vw}) = u_w(z_{wv})$ at every node,
\item Stability: every ghost component ($u_v$ constant) has degree $\operatorname{deg}(v) \ge 3$.
\end{enumerate}
\end{definition}

\begin{theorem}[Topological Ghost Component Bound]
\label{thm:ghost_components_bound}
\lean{GromovCompactness.ghost_components_bound}
\leanok
\uses{f06:sphere_tree}
In any stable sphere tree with $|V| > 1$, the number of ghost components $G$ satisfies:
\[
G \le P - 2,
\]
where $P$ is the number of non-constant (principal) components. Consequently, $|V| = P + G \le 2P - 2$.
\end{theorem}
\begin{proof}
\leanok
By the handshake lemma on finite trees: $\sum_{v \in V} \operatorname{deg}(v) = 2(|V| - 1) = 2(P + G - 1)$. Since ghost vertices satisfy $\operatorname{deg}(v) \ge 3$ and principal vertices satisfy $\operatorname{deg}(v) \ge 1$:
\[
3G + P \le \sum_{v \in V} \operatorname{deg}(v) = 2P + 2G - 2 \implies G \le P - 2.
\]
\end{proof}

% =========================================================================
\chapter{Uniform Local Analytic Estimates}

\section{Elliptic Estimates and Epsilon Regularity}

\begin{lemma}[Scalar Mean-Value Inequality]
\label{a01:mean_value}
\lean{GromovCompactness.Analysis.scalar_mean_value_inequality}
\uses{f00:library_audit}
Let $e \ge 0$ satisfy the subharmonic differential inequality $\Delta e \ge -b e^2$ on the Euclidean disk $D_r$. If $\int_{D_r} e < \eta = \frac{1}{2b}$, then:
\[
\sup_{D_{r/2}} e \le \frac{C}{r^2} \int_{D_r} e.
\]
\end{lemma}

\begin{lemma}[Bochner Energy-Density Inequality]
\label{a02:bochner}
\lean{GromovCompactness.Analysis.bochner_differential_inequality}
\uses{f01:target_geometry, f03:energy_measure}
The Riemannian energy density $e = \frac{1}{2}|du|^2$ of any $J$-holomorphic curve on a conformal chart satisfies the Bochner differential inequality:
\[
\Delta e \ge -b e^2 - c e,
\]
where the constants $b, c > 0$ depend only on the Riemann curvature of $(M, g)$ and $\|\nabla J\|_{C^1}$.
\end{lemma}

\begin{theorem}[$\varepsilon$-Regularity via Hofer's Lemma]
\label{a03:eps_regularity}
\lean{GromovCompactness.Analysis.epsilon_regularity}
\uses{a01:mean_value, a02:bochner, f00:library_audit}
There exist uniform constants $\varepsilon_{\mathrm{reg}} > 0$ and $C_{\mathrm{reg}} > 0$ such that for every $J$-holomorphic curve $u : D_r \to M$, if:
\[
E(u; D_r) < \varepsilon_{\mathrm{reg}},
\]
then:
\[
\sup_{D_{r/2}} |du|^2 \le \frac{C_{\mathrm{reg}}}{r^2} E(u; D_r).
\]
\end{theorem}

\begin{lemma}[Euclidean Elliptic and Sobolev Toolkit]
\label{a04:elliptic_toolkit}
\lean{GromovCompactness.Analysis.interior_cauchy_riemann_estimate}
\uses{f00:library_audit}
Interior $W^{k+1, p}$ estimates for the linear and perturbed Cauchy-Riemann operator $\delbar$: for $p > 2$,
\[
\|u\|_{W^{k+1, p}(D_{r/2})} \le C\big(\|\delbar u\|_{W^{k, p}(D_r)} + \|u\|_{L^p(D_r)}\big),
\]
along with the Morrey embedding $W^{1, p}(D_r) \hookrightarrow C^{0, 1 - 2/p}(D_r)$.
\end{lemma}

\begin{theorem}[Nonlinear Elliptic Bootstrap]
\label{a05:bootstrap}
\lean{GromovCompactness.Analysis.nonlinear_elliptic_bootstrap}
\uses{a04:elliptic_toolkit, f01:target_geometry, f04:smooth_local_convergence}
Every continuous weak solution $u \in W^{1, p}_{\mathrm{loc}}(\Omega, M)$ ($p > 2$) of the Cauchy-Riemann equation $\delbar_J u = 0$ is smooth ($C^\infty$).
\end{theorem}

\begin{theorem}[Local Compactness Under Gradient Bounds]
\label{a06:local_compactness}
\lean{GromovCompactness.Analysis.local_smooth_compactness}
\uses{a05:bootstrap, f04:smooth_local_convergence, f05:subsequence_bookkeeping, f00:library_audit}
Let $u_n : \Omega \to M$ be a sequence of $J$-holomorphic maps satisfying a uniform gradient bound $\sup_{K} |du_n| \le C_K < \infty$ on each compact subset $K \subset \Omega$. Then there exists a subsequence $u_{\sigma(n)}$ converging in $C^\infty_{\mathrm{loc}}(\Omega, M)$ to a $J$-holomorphic map $u : \Omega \to M$.
\end{theorem}

\section{Minimal Energy Gap and Singularity Removal}

\begin{theorem}[Minimal Energy Gap for Non-Constant Spheres]
\label{a07:sphere_energy_gap}
\lean{GromovCompactness.Analysis.sphere_energy_gap}
\uses{a03:eps_regularity, f01:target_geometry, f02:sphere_geometry, f03:energy_measure}
There exists a strictly positive constant $\varepsilon_{\mathrm{sph}} > 0$ depending only on $(M, \omega, J)$ such that every non-constant $J$-holomorphic sphere $u : \CP \to M$ satisfies:
\[
E(u) \ge \varepsilon_{\mathrm{sph}}.
\]
\end{theorem}

\begin{lemma}[Target Monotonicity]
\label{a08:monotonicity}
\lean{GromovCompactness.Analysis.target_monotonicity}
\uses{f01:target_geometry, f03:energy_measure, a04:elliptic_toolkit}
There exist $r_0 > 0$ and $c > 0$ such that for all $x \in M$ and $0 < r < r_0$, if a $J$-holomorphic curve $u$ intersects the metric ball $B_r(x)$ and has boundary outside $B_r(x)$, then:
\[
E\big(u; u^{-1}(B_r(x))\big) \ge c r^2.
\]
\end{lemma}

\begin{lemma}[Cylinder Localization]
\label{a09:cylinder_localization}
\lean{GromovCompactness.Analysis.cylinder_localization}
\uses{a03:eps_regularity, a08:monotonicity}
On a long conformal cylinder $[-T, T] \times S^1$, if the total energy is sufficiently small ($E \le \varepsilon_{\mathrm{loc}}$), the entire central segment $[-T/2, T/2] \times S^1$ maps into a single geodesic ball in $M$.
\end{lemma}

\begin{theorem}[Cylinder Exponential Decay]
\label{a10:cylinder_decay}
\lean{GromovCompactness.Analysis.cylinder_exponential_decay}
\uses{a09:cylinder_localization, a03:eps_regularity, f03:energy_measure}
On any conformal cylinder $[-T, T] \times S^1$ mapped into a coordinate chart, energy and loop diameter decay exponentially away from the ends:
\[
E(u; [t - 1, t + 1] \times S^1) \le C e^{-\delta(T - |t|)} E(u; [-T, T] \times S^1).
\]
\end{theorem}

\begin{theorem}[Sobolev Extension Across Punctures]
\label{a11:puncture_extension}
\lean{GromovCompactness.Analysis.puncture_sobolev_extension}
\uses{a09:cylinder_localization, a10:cylinder_decay, a03:eps_regularity, f03:energy_measure}
Any finite-energy $J$-holomorphic curve on a punctured disk $D^* = D_1 \setminus \{0\}$ extends to a continuous map in $W^{1, p}(D_1, M)$ ($p > 2$).
\end{theorem}

\begin{theorem}[Removal of Singularities]
\label{a12:removal_of_singularities}
\lean{GromovCompactness.Analysis.removal_of_singularities}
\uses{a11:puncture_extension, a05:bootstrap, f02:sphere_geometry}
Any $J$-holomorphic map $u : D_1 \setminus \{0\} \to M$ with finite energy $E(u) < \infty$ extends smoothly to a $J$-holomorphic map $\tilde{u} : D_1 \to M$.
\end{theorem}

% =========================================================================
\chapter{Concentration Compactness and Bubbling Engine}

\section{Measure Convergence and Moving Centers}

\begin{theorem}[Weak Compactness of Energy Measures]
\label{c01:weak_measures}
\lean{GromovCompactness.Compactness.weak_measure_compactness}
\uses{f03:energy_measure, f05:subsequence_bookkeeping}
Let $u_n : \CP \to M$ be a sequence of $J$-holomorphic spheres with uniformly bounded energy $E(u_n) \le B < \infty$. Then there exist a subsequence $u_{\sigma(n)}$ and a finite nonnegative Radon measure $\mu$ on $\CP$ such that $\mu_n = E(u_{\sigma(n)}; \cdot) \rightharpoonup \mu$ weakly.
\end{theorem}

\begin{theorem}[Finite Concentration Set and Limit Curve]
\label{c02:concentration_data}
\lean{GromovCompactness.Compactness.ConcentrationData}
\uses{c01:weak_measures, a03:eps_regularity, a06:local_compactness, a12:removal_of_singularities}
The concentration set $S := \{p \in \CP : \mu(\{p\}) \ge \varepsilon_{\mathrm{reg}}\}$ is finite with $|S| \le B / \varepsilon_{\mathrm{reg}}$. On $\CP \setminus S$, $u_{\sigma(n)} \to u_0$ in $C^\infty_{\mathrm{loc}}$, and $u_0$ extends smoothly across $S$ to a $J$-holomorphic curve $u_0 : \CP \to M$.
\end{theorem}

\begin{theorem}[Defect Measure Identity]
\label{c03:defect_identity}
\lean{GromovCompactness.Compactness.defect_measure_identity}
\uses{c02:concentration_data, f03:energy_measure}
The weak limit measure decomposes as:
\[
\mu = u_0^* \omega + \sum_{p \in S} m_p \delta_p,
\]
where each point mass satisfies $m_p \ge \varepsilon_{\mathrm{sph}}$, and the total energy is conserved: $\lim_{n \to \infty} E(u_{\sigma(n)}) = \mu(\CP) = E(u_0) + \sum_{p \in S} m_p$.
\end{theorem}

\section{The Bubbling Engine}

\begin{definition}[Universal Energy Threshold $h > 0$]
\label{b01:threshold_h}
\lean{GromovCompactness.Bubbling.TargetThresholds}
\uses{a03:eps_regularity, a07:sphere_energy_gap, a10:cylinder_decay}
Fix a target geometric constant $h > 0$ satisfying:
\[
0 < h \le \min\big(\tfrac{1}{2}\varepsilon_{\mathrm{reg}}, \, \varepsilon_{\mathrm{sph}}, \, \tfrac{1}{2}\varepsilon_{\mathrm{neck}}\big).
\]
\end{definition}

\begin{theorem}[Calibrated Rescaling Centers and Radii]
\label{b02:calibration}
\lean{GromovCompactness.Bubbling.CalibratedRescaling}
\uses{c03:defect_identity, b01:threshold_h, f02:sphere_geometry}
At each concentration point $p \in S$, there exist moving centers $a_n \to p$ and radii $\lambda_n \to 0$ satisfying the calibration condition $E(u_n; D_{\lambda_n}(a_n)) = m_p - h/2$.
\end{theorem}

\begin{theorem}[Vanishing Neck Energy and Diameter]
\label{b03:neck_vanishing}
\lean{GromovCompactness.Bubbling.neck_vanishing_quantified}
\uses{b02:calibration, a10:cylinder_decay, c03:defect_identity}
On annular neck regions $A(R\lambda_n, \delta; a_n)$, both the symplectic energy and target metric diameter vanish asymptotically as $n \to \infty, \delta \to 0, R \to \infty$.
\end{theorem}

\begin{theorem}[One-Level Mass Decomposition]
\label{b04:one_level_decomp}
\lean{GromovCompactness.Bubbling.OneLevelDecomposition}
\uses{b03:neck_vanishing, c02:concentration_data, c03:defect_identity, a12:removal_of_singularities}
Rescaled maps $v_n(z) := u_n(a_n + \lambda_n z)$ converge $C^\infty_{\mathrm{loc}}$ on $\C \setminus S_{\mathrm{child}}$ to a bubble $v : \CP \to M$, with $m_p = E(v) + \sum_{i} m_{p, i}$ where each child defect satisfies $m_{p, i} \ge h$.
\end{theorem}

\begin{theorem}[Ghost Branching Condition]
\label{b05:ghost_branching}
\lean{GromovCompactness.Bubbling.ghost_component_branches}
\uses{b04:one_level_decomp, b02:calibration}
If the limit bubble $v$ is a ghost component ($E(v) = 0$), the calibration condition guarantees at least 2 child defect points in $\C$, ensuring the vertex degree is at least 3 (accounting for the nodal attachment at $\infty$).
\end{theorem}

\begin{theorem}[Strict Energy Rank Decrease]
\label{b06:rank_decrease}
\lean{GromovCompactness.Bubbling.energy_rank_strict_decrease}
\uses{b04:one_level_decomp, b05:ghost_branching, b01:threshold_h}
Every child defect satisfies $m_{p, i} \le m_p - h$. Consequently, the discrete energy rank $\lfloor m / h \rfloor \in \N$ strictly decreases:
\[
\lfloor m_{p, i} / h \rfloor < \lfloor m_p / h \rfloor.
\]
\end{theorem}

\begin{theorem}[Local Defect Tree Resolution]
\label{b07:local_tree}
\lean{GromovCompactness.Bubbling.resolve_single_defect_tree}
\uses{b06:rank_decrease, f06:sphere_tree, f05:subsequence_bookkeeping}
By well-founded induction on the discrete energy rank $\lfloor m / h \rfloor \in \N$, every defect mass is resolved in finitely many bubbling steps into a finite rooted bubble tree.
\end{theorem}

% =========================================================================
\chapter{Global Tree Assembly and Master Compactness Theorem}

\section{Assembly and Stabilization}

\begin{theorem}[Prestable Tree Assembly]
\label{g01:global_assembly}
\lean{GromovCompactness.Assembly.PrestableLimitData}
\uses{b07:local_tree, c03:defect_identity, f02:sphere_geometry, f06:sphere_tree}
By composing and telescoping local rescaling frames, there exist global Möbius transformations $\phi_{v, n} \in \PSL$ for each vertex of the bubble tree such that $u_n \circ \phi_{v, n} \to u_v$ locally smoothly outside the nodal set.
\end{theorem}

\begin{theorem}[Root Stabilization]
\label{g02:root_stabilization}
\lean{GromovCompactness.Assembly.stabilize_prestable_tree}
\uses{g01:global_assembly, c02:concentration_data, f06:sphere_tree}
Contracting unstable ghost roots yields a fully stable sphere tree limit preserving all nodal matching conditions and energy conservation.
\end{theorem}

\section{The Master Compactness Theorem}

\begin{theorem}[Gromov Compactness Core in Genus Zero]
\label{g03:master_compactness}
\lean{GromovCompactness.gromov_compactness_core}
\uses{g02:root_stabilization, a07:sphere_energy_gap}
Every sequence of non-constant $J$-spheres $u_n : \CP \to M$ with uniformly bounded energy $E(u_n) \le B < \infty$ has a subsequence converging to a stable sphere tree limit.
\end{theorem}

\begin{theorem}[Universal Component and Vertex Bounds]
\label{g04:component_bounds}
\lean{GromovCompactness.Topology.component_count_bounds}
\uses{f06:sphere_tree, b01:threshold_h}
For any stable sphere tree limit with total energy bounded by $B$:
\begin{enumerate}
\item The number of principal components satisfies $P \le \lfloor B / h \rfloor$.
\item The total number of vertices satisfies $|V| \le \max\big(1, 2\lfloor B / h \rfloor + k - 2\big)$.
\end{enumerate}
\end{theorem}

\begin{theorem}[Master Integration Gate and Full Closure]
\label{g05:final_gate}
\lean{GromovCompactness.gromov_compactness_genus_zero}
\uses{g03:master_compactness, g04:component_bounds}
The complete master theorem uniting subsequence convergence, energy conservation, vanishing necks, and component bounds in Lean 4.
\end{theorem}

% =========================================================================
\chapter{Extensions}

\begin{theorem}[Preservation of Integral Homology Classes]
\label{e01:homology}
\lean{GromovCompactness.Extensions.integral_homology_preserved}
\uses{g05:final_gate, f01:target_geometry}
$[u_{\sigma(n)}] = \sum_{v \in V} [u_v]$ in $H_2(M; \Z)$ via 2-chain neck pinching.
\end{theorem}

\begin{theorem}[Component Stability Equals Finite Automorphisms]
\label{e02:automorphisms}
\lean{GromovCompactness.Extensions.stability_iff_finite_automorphisms}
\uses{f06:sphere_tree, a05:bootstrap}
The geometric stability condition $\operatorname{deg} \ge 3$ on ghost components is equivalent to finiteness of the tree automorphism group.
\end{theorem}

\begin{theorem}[Marked Point Compactness]
\label{e03:marked}
\lean{GromovCompactness.Extensions.marked_gromov_compactness}
\uses{g05:final_gate}
Compactness for sequences of $k$-marked $J$-holomorphic spheres via ratio-of-differences bubble insertion.
\end{theorem}

\begin{theorem}[Already Nodal Inputs]
\label{e04:nodal_input}
\lean{GromovCompactness.Extensions.nodal_input_gromov_compactness}
\uses{e03:marked, g04:component_bounds}
Extension of compactness to sequences of already-nodal genus 0 curves.
\end{theorem}

\begin{theorem}[Compact Hausdorff Moduli Space]
\label{e05:moduli_topology}
\lean{GromovCompactness.Extensions.moduli_space_compact_hausdorff}
\uses{e04:nodal_input, e02:automorphisms, e01:homology}
The compactified moduli space $\overline{\mathcal{M}}_{0, k}(M, A; J)$ is a compact Hausdorff topological space.
\end{theorem}
